\section{Case study and results}
\label{sec:case_study}

The case study evaluates the SAIL-SAC framework with P4 BalancedGate validation on the modified IEEE 33-bus feeder under network-stress conditions. Training uses 2023 load, PV, and price data over 100\,k steps with validation checkpoints every 2500 steps across four seasonal 30-day windows. Final evaluation uses the full 2024 year (365 days) to assess cross-year generalization. The study tests 16 random seeds to characterize seed-dependent failure modes, a critical dimension often underreported in DRL battery-control literature\cite{henderson2018deepRLmatters,agarwal2021statistical}.

\subsection{Multi-seed evaluation and success characterization}

Table~\ref{tab:seed_success_rate} reports seed-level success filtering. Of 16 tested seeds, 6 (37.5\%) satisfy the P4 morphology gates: mid-band SOC dwelling $>5\%$ and infeasible-action fraction $<30\%$. The remaining 10 seeds (62.5\%) exhibit degenerate policies with $>57\%$ infeasible actions and $<3\%$ mid-band dwelling, indicating that the learned policy often defaults to a battery-inactive or boundary-following mode rather than internalizing useful energy-shifting behavior. This failure rate is not merely random variation---it reflects a systemic challenge where SAC can fall into local optima in constrained sequential problems, producing acceptable-looking training curves without deployment-quality battery morphology.

\begin{table*}[!t]
\centering
\footnotesize
\setlength{\tabcolsep}{5pt}
\renewcommand{\arraystretch}{1.1}
\caption{Multi-seed success characterization. Gates: mid-band $>5\%$, infeasible $<30\%$.}
\label{tab:seed_success_rate}
\begin{tabular}{l>{\raggedright\arraybackslash}p{3.2cm}cc>{\raggedright\arraybackslash}p{6.2cm}}
\toprule
Category & Seeds & $n$ & Rate & Morphology \\
\midrule
Success & 42, 44, 81, 83, 84, 93 & 6 & 37.5\% & 8.5--49.1\% mid-band, 0--19.5\% infeasible. \\
Failure & 41, 43, 52, 62, 71, 82, 85, 91, 92, 94 & 10 & 62.5\% & 0.01--2.9\% mid-band, 57.6--99.9\% infeasible. \\
\bottomrule
\end{tabular}
\end{table*}

The six successful seeds display substantial morphological diversity. Seed 93 achieves the strongest mid-band performance (49.1\% dwelling, 0\% infeasible actions), while seeds 81, 83, and 84 exhibit more conservative mid-band residence (8.5--10.7\%) with differing infeasible-action fractions. Seed 42 provides an additional successful profile with 26.7\% mid-band dwelling and the strongest peak-price discharge response among the six seeds. This diversity indicates that successful policies can adopt different operational profiles while maintaining comparable economic performance. This observation is consistent with the multi-seed morphology scatter shown in Figure~\ref{fig:morphology_scatter}.

\subsection{Economic performance and cross-fidelity robustness}

Table~\ref{tab:six_seed_performance} reports the consolidated performance of the six successful seeds across four battery model fidelities. Mean objective cost is $11.788 \pm 0.120$\,M\textyen{} for the simple EBM baseline. The higher-fidelity evaluations are slightly lower in this logged dataset: 11.740\,M\textyen{} (Thevenin loss-only), 11.739\,M\textyen{} (Thevenin Rint-only), and 11.737\,M\textyen{} (Thevenin full PBM). Because objective accounting and battery-loss settlement differ across model fidelities, these values should be interpreted as small cross-fidelity variation rather than monotonic degradation. The simple-to-full mean difference is $-0.43\%$, supporting the engineering conclusion that policies trained on the simplified model retain comparable full-year performance under higher-fidelity battery settlement.

\begin{table*}[!t]
\centering
\footnotesize
\setlength{\tabcolsep}{6pt}
\renewcommand{\arraystretch}{1.12}
\caption{Six-seed consolidated performance across battery model fidelities (365-day cross-year evaluation, 2023 training $\to$ 2024 testing). Values are mean (standard deviation) over seeds 42, 44, 81, 83, 84, and 93.}
\label{tab:six_seed_performance}
\begin{tabular}{lcccc}
\toprule
Metric & Simple EBM & Thevenin Loss & Thevenin Rint & Thevenin Full PBM \\
\midrule
Objective cost (M\textyen{}) & 11.788 (0.120) & 11.740 (0.104) & 11.739 (0.106) & 11.737 (0.103) \\
Mid-band SOC dwell (\%) & 20.4 (15.6) & 19.0 (14.7) & 19.8 (14.8) & 18.6 (14.6) \\
Peak-price discharge (\%) & 24.7 (26.3) & 38.7 (29.7) & 39.6 (30.1) & 38.6 (29.5) \\
Valley-price charge (\%) & 27.7 (9.6) & 36.7 (8.8) & 36.7 (8.4) & 37.2 (8.9) \\
Infeasible actions (\%) & 6.1 (7.3) & 7.7 (10.3) & 7.7 (10.4) & 7.8 (10.4) \\
Power-flow failures & 0 (0) & 0 (0) & 0 (0) & 0 (0) \\
\bottomrule
\end{tabular}
\end{table*}

\subsection{Intervention diagnostics and price-response behavior}

Table~\ref{tab:shield_dependence} reports intervention-related diagnostics for the six successful seeds under the simple EBM full-year evaluation. In this 16-seed P4 dataset, \texttt{shield\_enabled=0} and all logged shield-activation fields are zero. Therefore, the current evidence does not support claims of residual deployed shield activation. The relevant nonzero diagnostic is the inventory-teacher action gap: the six successful seeds have mean absolute gap $0.301 \pm 0.087$, indicating that the learned actions remain separated from the rule/teacher reference even when the deployment shield is not active.

\begin{table}[!t]
\centering
\footnotesize
\setlength{\tabcolsep}{5pt}
\renewcommand{\arraystretch}{1.12}
\caption{Intervention and price-response diagnostics for six successful seeds (365-day evaluation, simple EBM). Values are mean (standard deviation).}
\label{tab:shield_dependence}
\begin{tabular}{lc}
\toprule
Metric & Mean (Std) \\
\midrule
Mean absolute shield delta & 0.000 (0.000) \\
Material shield activation & 0.000 (0.000) \\
Strong shield activation & 0.000 (0.000) \\
Inventory-teacher action gap & 0.301 (0.087) \\
Peak-price discharge fraction & 0.247 (0.263) \\
Valley-price charge fraction & 0.277 (0.096) \\
\bottomrule
\end{tabular}
\end{table}

Figure~\ref{fig:shield_dependence} visualizes the same diagnostics. Panel (a) compares the per-seed inventory-teacher gap against the shield delta, making clear that the former is nonzero while the latter is zero in the current full-year audit. Panel (b) shows that any, material, and strong shield activation are all zero. Panel (c) summarizes SOC morphology, price-response, infeasible dwell, shield delta, and teacher gap with inter-seed variability.

\begin{figure*}[!t]
\centering
\includegraphics[width=0.95\textwidth]{fig1_shield_dependence_elite.pdf}
\caption{Intervention diagnostics for the six successful seeds. (a) Per-seed mean absolute inventory-teacher gap and logged shield delta. (b) Shield-activation fractions; all are zero in the current P4 full-year dataset. (c) Mean and standard deviation of morphology, price-response, infeasible-action, shield-delta, and teacher-gap metrics.}
\label{fig:shield_dependence}
\end{figure*}

The strict six-seed audit also changes the price-response interpretation. Mean peak-price discharge fraction is 24.7\%, and mean valley-price charge fraction is 27.7\%, with large inter-seed dispersion. Seed 42 exhibits strong peak discharge (75.1\%), but several other successful seeds satisfy the morphology gates with much weaker price-directional behavior. Thus, the data support a more cautious conclusion: P4 selection identifies policies with acceptable SOC morphology and infeasible-action suppression, but price-responsive arbitrage remains heterogeneous and should not be overstated.

\subsection{Protocol ablation: P1--P4}

Table~\ref{tab:p1_p4_tradeoff} compares the four validation variants on seed 42, the most extensively analyzed seed in the study. These values summarize protocol-window diagnostics and are separate from the full-year six-seed P4 audit above. P1 (strict) is excessively conservative: positive value recovery requires willingness to accept some intervention, which P1 penalizes so heavily that no checkpoint passes its gates. P2 (no-strict) achieves the highest raw recovery (0.598) but at substantially elevated intervention dependence (35.2\% strong activation). P3 (coarse checkpointing) collapses entirely under infrequent validation: 79.2\% strong shield activation indicates the selected checkpoint relies almost entirely on the safety layer for constraint satisfaction. P4 BalancedGate occupies the most operationally useful position in this protocol diagnostic: positive recovery (0.456), meaningful mid-band dwelling (30.1\%), strong peak-discharge behavior (76.8\%), and lower strong shield activation (22.9\%) than P2 and P3.

\begin{table*}[!t]
\centering
\footnotesize
\setlength{\tabcolsep}{6pt}
\renewcommand{\arraystretch}{1.12}
\caption{P1--P4 validation protocol ablation (seed 42, 30-day cross-year windows). Raw reference-normalized recovery used for diagnostic comparison.}
\label{tab:p1_p4_tradeoff}
\begin{tabular}{lccccc p{4.2cm}}
\toprule
Variant & Recovery & Mid-band & Peak disch. & Valley chg. & Strong shield & Interpretation \\
\midrule
P1 strict & 0.025 & 0.075 & 0.309 & 0.341 & 0.105 & Overly conservative; weak value and morphology. \\
P2 no-strict & 0.598 & 0.209 & 0.604 & 0.481 & 0.352 & Higher value; elevated intervention dependence. \\
P3 coarse & 0.166 & 0.101 & 0.312 & 0.412 & 0.792 & Collapse under coarse schedule. \\
P4 balanced & 0.456 & 0.301 & 0.768 & 0.498 & 0.229 & Balanced; positive value with managed shield. \\
\bottomrule
\end{tabular}
\end{table*}

If only objective recovery were reported, P2 would dominate. If only conservative shield behavior were emphasized, P1 would appear preferable. The combined value--morphology view identifies P4 as the engineering compromise in this seed-42 diagnostic: it improves storage behavior while preserving a transparent record of residual intervention dependence, without overstating autonomous policy capability.

Figure~\ref{fig:protocol_tradeoff} maps the same ablation in the recovery--midband plane, with marker color encoding peak-discharge fraction and marker size scaling with strong shield activation. P4 occupies the most balanced quadrant, combining non-trivial recovery with higher mid-band dwell and peak-discharge fraction than P1 and P3, while avoiding the severe intervention dependence of P3 (79.2\%).

\begin{figure}[!t]
\centering
\includegraphics[width=\columnwidth]{fig_protocol_tradeoff.pdf}
\caption{P1--P4 value--morphology tradeoff. (a) Raw reference-normalized recovery versus SOC mid-band dwell; marker color denotes peak-price discharge fraction and marker size increases with strong shield activation. (b) Strong shield activation per variant.}
\label{fig:protocol_tradeoff}
\end{figure}

\subsection{Baseline comparison}
\label{sec:external_baselines}

Table~\ref{tab:external_baseline_comparison} compares P4 BalancedGate against two within-family baselines: SAC with the same shield but cost-only checkpoint selection, and vanilla SAC without corrective shielding. The baseline columns are diagnostic rather than paired statistical comparisons because the seed sets and shield settings differ. The cost-only SAC+shield baseline treats the shield as invisible infrastructure and selects by scalar reward. The vanilla SAC baseline removes corrective shielding, retaining only physical battery clipping.

\begin{table}[!t]
\centering
\scriptsize
\setlength{\tabcolsep}{3.2pt}
\renewcommand{\arraystretch}{1.12}
\caption{Baseline comparison on the IEEE 33-bus network-stress case (six successful seeds for P4; seeds 42/52/62 for cost-only+shield; seeds 42/52 for vanilla SAC).}
\label{tab:external_baseline_comparison}
\begin{tabular}{lccc}
\toprule
Metric & P4 & Cost+shield & Vanilla SAC \\
\midrule
Mid-band SOC dwell & \textbf{20.4\%} & 0.0\% & 0.4\% \\
Peak disch.\ fraction & 24.7\% & \textbf{44.1\%} & 25.9\% \\
Material shield activ. & \textbf{0.0\%} & 65.9\% & --- \\
Infeasible dwell & 6.1\% & \textbf{0.03\%} & 23.6\% \\
Power-flow failures & 0 & 0 & 0 \\
\bottomrule
\end{tabular}
\end{table}

Three findings emerge from this diagnostic comparison. First, P4 BalancedGate substantially improves mid-band SOC dwelling relative to cost-only and vanilla selection, which is the central morphology objective. Second, cost-only+shield achieves near-zero infeasible dwell but does so with high material shield activation in the reported seed subset, illustrating why shield metrics must be visible rather than hidden. Third, vanilla SAC without corrective shielding produces much higher infeasible-action dwell, confirming that unconstrained SAC is not a reliable deployment controller in this network-stress setting.

\subsection{Operational interpretation}

The results are limited but the limitations are quantified and visible. The six successful seeds achieve positive economic performance ($11.788$\,M\textyen{} mean simple-model cost), zero power-flow failures, 20.4\% mean mid-band dwell, and 6.1\% mean infeasible-action dwell. Their morphological range---8.5\% to 49.1\% mid-band dwelling---demonstrates that the P4 protocol can identify operationally meaningful policies across diverse behavioral profiles. At the same time, the current full-year audit shows zero deployed shield activation and a nonzero inventory-teacher gap, so the correct interpretation is not residual shield activation but incomplete internalization relative to the teacher reference. The LP oracle normalization reference assumes perfect foresight; the resulting recovery index is an engineering diagnostic rather than an optimality certificate.

The multi-seed framework Figure~\ref{fig:morphology_scatter} and the cross-fidelity Figure~\ref{fig:cross_fidelity_heatmap} together provide the two core visual evidences of the paper: the former shows that morphological success is non-trivial to achieve (37.5\% rate) and that successful seeds occupy a geometrically distinct region; the latter shows that cross-fidelity variation is small across the six successful seeds, validating policy transfer robustness without claiming monotonic degradation.

\begin{figure*}[!t]
\centering
\includegraphics[width=0.95\textwidth]{fig2_cross_fidelity_elite.pdf}
\caption{Cross-fidelity performance for six successful seeds. (a) Objective cost across simple EBM, Thevenin loss-only, Thevenin Rint-only, and Thevenin full PBM, with per-seed traces and mean $\pm 1\sigma$. (b) Morphology score (mid-band\% $-$ infeasible\%) from simple EBM to Thevenin full PBM for each seed. All values use 365-day full-year evaluation.}
\label{fig:cross_fidelity_heatmap}
\end{figure*}

\begin{figure}[!t]
\centering
\includegraphics[width=\columnwidth]{fig3_morphology_density_elite.pdf}
\caption{Morphology scatter for all 16 evaluated seeds. The success region is defined by mid-band dwell $>5\%$ and infeasible-action dwell $<30\%$, enclosing 6/16 seeds. Failure seeds cluster at high infeasible dwell and near-zero mid-band residence. Successful points are colored by objective cost, and labels identify the six successful seeds.}
\label{fig:morphology_scatter}
\end{figure}
